\documentclass[11pt]{article}
\usepackage{icmlko}
\begin{document}

\title{어느 행을 더 모으나 --- 그리고 만화는 받기만 하고 주지 않는다}
\author{월드모델 실험실}
\date{노트 508 · 2026년 8월 3일}
\maketitle

\begin{abstract}
노트 506 이 ``행을 두 배로 하면 $+0.0245$''를 냈다. 다음 물음은
\textbf{어느 도메인의 행인가}이다. 도메인 하나씩 학습 행을 절반으로 줄여
판 손해를 쟀다(짝 씨앗 $6$ · 나머지는 그대로).

\begin{center}
\begin{tabular}{lrrrr}
\hline
도메인 & 뺀 행 & 판 손해 & 부호 & 행 만 개당 \\
\hline
\textbf{모바일} & $779$ & $-0.0048$ & $\mathbf{6/6}$ & $\mathbf{-0.0610}$ \\
애니 & $1{,}389$ & $-0.0048$ & $5/6$ & $-0.0344$ \\
세계애니 & $1{,}324$ & $-0.0036$ & $6/6$ & $-0.0269$ \\
펀딩 & $1{,}179$ & $-0.0031$ & $5/6$ & $-0.0266$ \\
웹툰 & $1{,}471$ & $-0.0038$ & $4/6$ & $-0.0259$ \\
게임 & $129$ & $-0.0002$ & $3/6$ & $-0.0188$ \\
\hline
\textbf{만화} & $\mathbf{2{,}823}$ & $\mathbf{+0.0026}$ & $1/6$ & $\mathbf{+0.0091}$ \\
\hline
\end{tabular}
\end{center}

사전 등록(``행당 값어치가 다섯 배 넘게 갈린다'')은 \textbf{통과} ---
$-0.0610$ 에서 $+0.0091$ 까지 $6.7$ 배다. \textbf{수집 우선순위가 의미
있다.}

그리고 하나가 부호까지 다르다 --- \textbf{만화는 학습 행 $2{,}823$
개(전체 학습의 $15\%$)를 빼면 판이 $+0.0026$ 오른다.}
\end{abstract}

\begin{figure}[t]
\centering
\includegraphics[width=\textwidth]{figs/f1.pdf}
\caption{\textbf{가} --- 행당 값어치. \textbf{나} --- 받는 쪽과 주는 쪽.
\textbf{다} --- 작은 도메인은 못 잰다.}
\end{figure}

\section{받는 쪽과 주는 쪽}

노트 505 가 도메인마다 ``챔피언 $\div$ 자기 천장''을 쟀다 --- 그것이
\textbf{얼마나 받는지}다. 이 노트가 잰 행당 값어치는 \textbf{얼마나
주는지}다. 둘을 함께 놓으면 만화가 혼자 떨어져 나온다.

\begin{center}
\begin{tabular}{lrr}
\hline
도메인 & 받는다(천장 대비) & 준다(행 만 개당) \\
\hline
\textbf{만화} & $\mathbf{2.00}$ & $\mathbf{+0.0091}$ (해롭다) \\
펀딩 & $1.72$ & $-0.0266$ \\
모바일 & $1.11$ & $-0.0610$ \\
세계애니 & $1.07$ & $-0.0269$ \\
애니 & $1.03$ & $-0.0344$ \\
게임 & $1.02$ & $-0.0188$ \\
웹툰 & $1.01$ & $-0.0259$ \\
\hline
\end{tabular}
\end{center}

\textbf{만화는 제일 많이 받고(자기 천장의 두 배) 유일하게 안 준다.} 학습
행 $5{,}646$ 개로 열한 도메인 중 제일 많은데(전체의 $30\%$) 한계 기여가
음수다.

펀딩도 많이 받지만($1.72$) 준다($-0.0266$). \textbf{받는 것과 주는 것은
따로다.}

\section{작은 도메인은 이 자로 못 잰다}

도서 $40$ 행 · 시장팝업 $61$ · 아이돌 $61$ · 팝업 $9$ 행만 빠져서 행당
수치가 튄다(팝업 $+0.6170$ 은 \textbf{$9$ 행짜리}다).

\begin{quote}
\textbf{조항} --- \textbf{행당 값어치는 뺀 행이 $100$ 개 이상일 때만
적는다.} 절반으로 줄여도 $\texttt{MIN\_TRAIN}$ 문턱에 걸려 거의 안 빠지는
도메인이 넷이다.
\end{quote}

\section{다음 --- 만화를 관문에}

\textbf{만화 학습 행 절반}을 자 다섯으로 태우고 있다(짝 씨앗 $20$).
통과하면 노트 507 조항대로 씨앗 $40$ 으로 확인한다.

기대는 조심스럽다 --- $+0.0026$ 은 노트 507 이 방금 만든 문턱
($+0.002$) 바로 위다. \textbf{세계애니 \texttt{grp} 이 $+0.0011$ 에서
죽은 것을 방금 봤다.}

\section*{남은 것}

\begin{itemize}
\item \textbf{모바일 행을 늘리면} --- 행당 $-0.0610$ 로 $1$ 등이고 지금
학습 행이 $1{,}558$ 개로 적다. \textbf{두 배로 늘리면 $+0.0095$}(이
자의 외삽).
\item \textbf{만화가 왜 안 주나} --- 노트 505 는 만화 자기 천장이
$+0.1679$ 로 낮다고 했다. \textbf{라벨이 시끄러운가, 아니면 다른
도메인과 눈금이 다른가.}
\item \textbf{행당 값어치와 판 무게는 다르다} --- 웹툰은 판 무게 $20.7\%$
로 $1$ 등인데 행당 값어치는 $5$ 등이다.
\end{itemize}

\end{document}
